\FloatBarrier
\subsection{每个保存目标时刻的位置误差}
\label{sec:point-error-horizons}
ES 评价整个预测分布，不是分布均值到真实位置的距离。表~\ref{tab:point-error-horizons} 将此前 Full/GMM/dt300 对照与惯性参考的这个距离单独列出，不构成全部28配置的新排行。使用原保存评分圆盘的中心 $c_{m,b,s,k}$（512粒子的均值，或确定性参考位置）和同一原目标 $y_{b,k}$，两者均在原起点局部评分坐标系中。不需打开预测数组、重新计算粒子评分、插值目标或模拟。

对模型 $m$，先在每个记录哈希块内平均其 $S_m$ 个原预测种子的误差，再等权平均全部46块。三个概率配置的 $S_m=5$，惯性参考的 $S_m=1$，不将一条参考计成五份：
\begin{equation}
\bar e_{m,k}=\frac{1}{46}\sum_{b=1}^{46}
\frac{1}{S_m}\sum_{s=1}^{S_m}\lVert c_{m,b,s,k}-y_{b,k}\rVert_2.
\label{eq:point-error-horizon-mean}
\end{equation}

\begin{table}[htbp]\centering\small
\caption{原保存均值位置的平均位移误差，单位米，越低越好。四列是原名义1/5/15/30分钟目标；保存时钟实际范围49--67、294--305、889--910、1785--1814秒，原30秒容差不变。份数是保存预测，不是独立目标或人员。仅描述数值，不做分时域推断。}\label{tab:point-error-horizons}
\begin{tabular}{@{}lrrrrr@{}}\toprule
模型 & 份数 & 1 分钟 & 5 分钟 & 15 分钟 & 30 分钟\\\midrule
Full & 230 & 73.77 & 332.31 & 812.26 & 1326.09\\
GMM & 230 & 71.92 & 323.79 & 792.99 & 1299.11\\
dt300 & 230 & 75.88 & 340.10 & 825.76 & 1335.73\\
惯性参考 & 46 & 31.17 & 237.93 & 934.03 & 2095.78\\
\bottomrule\end{tabular}\end{table}

全部736份所选保存结果对应同一184个目标位置。逐份检查四时刻误差的算术均值及最后一项，在浮点舍入精度内分别复现原保存 ADE 和 FDE，最大差异小于 $10^{-12}$ 米。先跨种子平均位置再取欧氏距离会得到另一估计量，本表未这样处理。

惯性参考在1、5分钟的平均点误差更低，Full 在15、30分钟更低。GMM 在四时域的描述性均值均略低于 Full；相反，dt300 虽有更低的原汇总 ES，四时域均值位置误差却均大于 Full。因此，分布评分改善不自动等于均值路径位置改善。dt300 指训练观测间隔，不是积分步长。事后均值不证明统计优越性、独立人员泛化或部署精度；原最终推断、输出核验及审阅仍待完成。
\FloatBarrier
